% Source: sources/returns/S6_04_MAPS_MANIFOLDS/sections/manifolds.tex:3579–3675
\paragraph{Corollary D.2.17 (Proved: two further original additive profiles).}
The complete formulas and oriented diffeomorphism classification
(D.T1)--(D.T3), for every free smooth-fibre log, also hold for
\[
\begin{array}{c|c|l|r|r}
 \mathrm{OS}&\mathrm{profile}&\text{ordered fibres}&p&b\\\hline
 46&II&I_3^*,II,I_1^\bullet&4&2\\
 50&IV&I_1^*,IV,I_1^\bullet&12&1.
\end{array}
\]
Use $P=pG$, $Q=G$, the simple linearization zero shown,
and the original fillings at both additive fibres.
\begin{proof}
The section heights and component residues are the same at
the star fibre as in D.2.16. At $II$, both sections are
narrow. At $IV$, $P$ is narrow and $Q$ has residue two,
a generator of its order-three component group.
The mixed height is again one. The same zero-multidegree
argument gives the required degree-one linearization, with
no zero at either additive fibre.

For $II$, contract both section logarithms to zero. The
marked original filling is $E_{II}\times T^2$, where
$E_{II}$ retracts onto its cuspidal rational curve,
topologically $S^2$. For $IV$, divide $3Q$ in the additive
identity component as in D.2.16. The resulting section
$Q_0$ has order three and cycles all three components of
$IV$. Its central curve is three spheres meeting at one
point, so $H_2(E_{IV})=\mathbb Z^3$ and $H_1(E_{IV})=0$.
The marked fill is
$S^1_\delta\times T(t_{Q_0}:E_{IV}\to E_{IV})$.
The local logarithm, translation extension and boundary
rigidification arguments apply verbatim. In either case
there is one component orbit, and hence
\[
 \pi_1(U)=\mathbb Z^2,\qquad
 (H_0,H_1,H_2,H_3,H_4)(U)
       =(\mathbb Z,\mathbb Z^2,\mathbb Z^2,\mathbb Z^2,\mathbb Z).
\]

The two additive monodromies in the respective global
markings are
\[
 A_{II}=\begin{pmatrix}
 1&1&0&0\\-1&0&0&0\\0&0&1&0\\0&0&0&1
 \end{pmatrix},\qquad
 A_{IV}=\begin{pmatrix}
 -1&1&0&-2\\-1&0&0&0\\-24&0&1&-16\\0&0&0&1
 \end{pmatrix}.
\]
Both $A-I$ and $\bigwedge^2A-I$ have only unit nonzero
Smith entries, of ranks two and four. Thus their boundaries
have $H_1=\mathbb Z^3$, $H_2=\mathbb Z^4$. The two circle
actions sweep the origin disc and give the complete
primitive normal-torus kernel, with fixed vectors
\[
 \delta,\gamma\quad\text{at }II;\qquad
 \delta,(-2,2,0,3)\quad\text{at }IV.
\]
The second vector has minimal positive $\gamma$-coordinate
three, exactly the period of the order-three translation.
The relative sequence, with $H_2(U,\partial U)=\mathbb Z$
mapping onto the meridian kernel, shows that the boundary
specialization is onto $H_2(U)$. Thus these additional
original fillings contribute no free defect or torsion.

There are now four punctures including the regular log.
The graph differential has domain $\mathbb Z^{12}$ and
rank three with primitive image. The six unchanged normal
columns, together with $j_0(e_1),j_0(e_2),j_0(\gamma)$,
form a primitive basis of its rank-nine kernel. Its integer
left inverse has the same residual projection
$(x_1,x_2,x_4)$. Exterior coinvariants are still detected
primitively by $p e_1^*e_2^*-\delta^*\gamma^*$.
Consequently the entire residual normal matrix is (D.T10),
with the same free star defect $\mathbb Z^b$.
The degree-one Mayer--Vietoris source and target now have
rank fourteen; their kernel is zero for $n\ne0$ and is
$\mathbb Z$ for $n=0$. The full homology therefore follows
by the identical unrestricted Smith proof and duality.

Finally the global $\delta$ action remains free on both
original additive fillings. Its unique fixed curve is
still the null-homologous cusp curve for $n=\pm1$.
The free $H_2$ remains represented by the star neighbourhood,
where the canonical line is trivial. Thus the projectably
spin surgery, Smale--Barden orbit classification, repeated
Galaz-Garc\'ia--Reiser Theorem A and Zhubr cancellation in
D.2.16 apply without modification and give (D.T3).

The independent program
\texttt{manifolds\_star\_additive.py} records the full
constant integral certificates and 1,462 full $10\times11$
Smith checks, including 194 cases with $n=0$. The proof is
unrestricted in the affine vector and multiplicity. No
priority claim is made.
\end{proof}
